% SPDX-License-Identifier: GPL-3.0-or-later
% Copyright (C) 2026 Alessandro Greco
\chapter{Предисловие}

Tor ежедневно используют миллионы людей, однако о нём почти всегда говорят через призму его реальных или предполагаемых применений: обход цензуры, защита источников, \enquote{тёмная сеть}. Эта книга подходит к теме иначе: она объясняет, как работает Tor. Мы начинаем с байтов, которыми обмениваются два реле, а затем, глава за главой, переходим к архитектурным решениям, атакам, реальным случаям и вопросам, которые по-прежнему остаются открытыми. И читателю не предлагается принимать изложенное на веру: каждый механизм, описанный в технической части, создаётся и проверяется в лаборатории, в частной сети Tor, которую можно запустить на обычном компьютере.

\section*{Происхождение, метод и приглашение к проверке}

В основу этой книги легли учебные заметки и документы, которые я составлял на протяжении многих лет, начиная с бакалаврской выпускной работы. Чтобы превратить эти материалы в связное техническое изложение, я руководил использованием инструментов искусственного интеллекта. Они внесли существенный вклад в разработку объяснений, аргументацию и написание текста, а также в создание иллюстраций, упражнений, кода и лабораторных работ. ИИ не просто редактировал мой текст и не заменяет независимую экспертную проверку.

Несмотря на технические проверки и мою собственную редактуру, в книге могут оставаться ошибки, упрощения и устаревшие сведения. Прошу читателей сообщать о них и предлагать исправления через репозиторий проекта. Я решил распространять книгу бесплатно через официальные каналы проекта, а её исходные материалы --- под лицензией GNU GPL-3.0-or-later, следуя принципу \enquote{diffusione dei saperi senza fondare poteri}: распространять знания, не создавая центров власти. Условия повторного использования, в том числе коммерческого, определяет лицензия.

\section*{Для кого эта книга}

Книга предназначена для студентов, изучающих информатику и информационную безопасность, специалистов, которые сталкиваются с Tor в работе, людей, уже управляющих реле или планирующих это делать, а также исследователей. От читателя требуется базовое понимание сетей (TCP/IP и общая идея TLS), знакомство с командной строкой и, для лабораторных работ, Python. Необходимые криптографические основы вводятся в \cref{cap:05-crittografia}, а дополнительные сведения приведены в \cref{cap:B-crittografia}.

\section*{Как читать эту книгу}

Книга состоит из пяти частей. Их можно читать последовательно или выбрать один из следующих маршрутов:
\begin{description}
  \item[Исторический и концептуальный маршрут] Часть I (главы 1--4) и часть V (главы 29--31): что означает анонимность, откуда появился Tor, от каких противников он защищает и в каком направлении развивается.
  \item[Протокольный маршрут] Часть II (главы 5--16), которую лучше читать по порядку вместе с лабораторными работами: это техническое ядро книги.
  \item[Практический маршрут] Часть III (главы 17--25): настройка tor, наблюдение за его трафиком, публикация onion-сервиса, эксплуатация реле, использование Tor Browser, измерение сети и понимание правовых аспектов. \Cref{cap:19-rete-privata} является обязательной основой для всех лабораторных работ.
  \item[Маршрут по безопасности] Часть IV (главы 26--28), перед которой следует прочитать \cref{cap:04-modello-minaccia}: атаки рассматриваются через исследования, реальные случаи и эксплуатационные ошибки.
\end{description}

\section*{Условные обозначения}

Каждая глава начинается с учебных целей и строки \emph{Последняя проверка}, где указаны дата, версии программного обеспечения и спецификаций, по которым проверялось содержание. Tor меняется, поэтому читателю важно понимать, какую версию Tor описывает текст. Во всей книге \emph{Tor} обозначает сеть и проект, \emph{tor} --- программу на C, а \emph{Arti} --- реализацию на Rust. Ссылки на спецификации записываются в форме \spec{document}{section}; предложения об изменении протокола идентифицируются по номеру. Принятые в русском издании технические термины вместе с их английскими эквивалентами собраны в \cref{cap:A-glossario}.

Врезки отделяют основной текст от отступлений: \emph{подробных разборов}, \emph{исторических справок} и \emph{предупреждений}. Сложность упражнений обозначается звёздами: одна --- на понимание, две --- на применение, три --- лабораторная задача или небольшой исследовательский проект; решения находятся в \cref{cap:D-soluzioni}. Лабораторные работы пронумерованы от 1 до 20. В части II они постепенно создают небольшой клиент Tor на Python (\texttt{torlab}); код и лабораторная сеть сопровождают книгу, а автоматизированные проверки можно повторно запускать после обновлений. Все эти ресурсы --- пакет \texttt{torlab}, двадцать лабораторных работ, тестовая сеть Docker и инструменты проверки --- доступны в репозитории книги: \repolibro.

\section*{Этическая оговорка}

Все эксперименты в этой книге выполняются в частной лабораторной сети Tor либо используют действующую сеть только как обычный клиент, без наблюдения за чужим трафиком и без его изменения. Именно такой принцип рекомендует исследователям Tor Research Safety Board; он также имеет правовое значение (\cref{cap:25-legale}). Главы об атаках написаны с аналитической и защитной точки зрения: они объясняют, что происходило и как от этого защищаться, а не как атаковать конкретного человека.

\section*{Благодарности}

Эта книга опирается на работу Tor Project и его сообщества. Открытые спецификации, исходный код \texttt{tor} и Arti, предложения об изменении протокола и документация лежат в основе всех последующих технических утверждений. Благодаря их публичной доступности утверждения можно сверять с первоисточниками, а не принимать на веру. Благодарности заслуживают и разработчики свободных инструментов, с помощью которых создавалась книга и тестировались её лабораторные работы: \LaTeX{}, Python, Docker, Chutney, Wireshark и многие другие проекты, упомянутые в тексте.
